\documentclass[12pt, a4paper]{article}

\usepackage[margin=1in]{geometry}
\usepackage{booktabs}
\usepackage{array}
\usepackage{enumitem}
\usepackage{hyperref}
\usepackage{microtype}
\usepackage{parskip}
\usepackage{xcolor}
\usepackage{titlesec}

\hypersetup{
    colorlinks=true,
    linkcolor=blue,
    urlcolor=blue
}

\titleformat{\section}{\large\bfseries}{}{0em}{}[\titlerule]
\titlespacing{\section}{0pt}{12pt}{6pt}

\title{\textbf{Guidelines for Course Project}}
\date{}

\begin{document}

\maketitle
\thispagestyle{empty}

To finalize the course project topic, select 3--4 recent papers from high-quality conferences
and journals to set the baseline. One can compare the results from different papers or try
different controllers studied in class. Before finalizing the topic, you must get my approval.

The project carries \textbf{30\% weightage}, which will be divided into the following stages:

\begin{itemize}[leftmargin=1.5em]
    \item \textbf{Stage 1:} Finalization of the problem statement and its submission as a report. At
    this stage, submit a preliminary report containing a description of the system and
    the problem statement.

    \item \textbf{Stage 2:} Understanding the methodology and proposing a new methodology along
    with analysis and submitting it as a report.

    \item \textbf{Stage 3:} Submission of the final project report and working code for implementation.

    \item \textbf{Stage 4:} Submission of the presentation.

    \item \textbf{Stage 5:} Final presentation with Q\,\&\,A.
\end{itemize}

\vspace{1em}

\begin{center}
\begin{tabular}{@{}lcc@{}}
\toprule
\textbf{Stage} & \textbf{Weightage} & \textbf{Deadline} \\
\midrule
Stage 1 & 3\%  & 10th March 2026 \\
Stage 2 & 4\%  & 31st March 2026 \\
Stage 3 & 9\%  & 5:30\,pm, 30th April 2026 (Thursday) \\
Stage 4 & 4\%  & 5:30\,pm, 1st May 2026 (Friday) \\
Stage 5 & 10\% & 2nd May 2026 (Saturday), 9:30\,am -- 1:00\,pm \\
\bottomrule
\end{tabular}
\end{center}

\vspace{0.5em}

One can earn \textbf{extra points} by proposing a novel idea that could lead to a
research article. Marks in the range $[5,\,10]$ will be added to the final total of
the course out of 100.

\newpage

\section*{Guidelines for Report Preparation}

\begin{enumerate}[leftmargin=1.5em]
    \item Prepare the report in \LaTeX{}. Use \texttt{documentclass} \texttt{article} or \texttt{report}. Single-column format.

    \item Equations should be typed, not copied from sources.

    \item Prepare figures using any tool you prefer, but do not copy-paste the figure unless it
    is very complicated. Give a proper citation of the figure if you are using the existing
    one from the literature.

    \item The plots generated from the code should be included in the report and PPT with
    proper thickness of line, font size of 14 for axes, and label so that figures are properly
    visible.

    \item You can prepare a PPT in Word or \LaTeX{}.

    \item Include a list of references and cite them properly.
\end{enumerate}

The presentation time is \textbf{15 min}, and \textbf{5 min} for Q\,\&\,A.

All of you are required to be present for the whole duration.

\vspace{1em}
\noindent\textit{Note: I will keep updating this PDF if something else needs to be added.}

\end{document}
